\documentclass{situationist-preprint}

\title{One Cloud: The Problem of the Many Resolved}
\msid{MET-05-a}
\msfamily{MET-05}
\msdiscipline{Metaphysics}
\mstype{Principal manuscript}
\mskeywords{problem of the many; clouds; supervaluation; epistemicism; vagueness; counting; maximality}

\begin{document}
\maketitle

\begin{abstract}
A cloud has no sharp edge: at its boundary, water droplets thin out gradually, and many aggregates of droplets, differing only at the edge, have an equal claim to be the cloud. If each is a cloud, there are many clouds where we would count one; if none is, there are no clouds; and choosing one seems arbitrary. We argue that there is exactly one cloud. The many candidates are not each a cloud: \enquote{cloud} is a maximal sortal, so that no cloud is a large proper part of another, and every admissible precisification of the term counts exactly one cloud in the region. \enquote{There is exactly one cloud there} is therefore true on every precisification, although no candidate is the cloud on all of them. We reconcile this supervaluationist treatment with the epistemic account of vagueness defended in MET-18: supervaluation describes what is determinately, that is knowably, true, while on the epistemic account exactly one candidate is in fact the referent of \enquote{the cloud}, and which one cannot be known.
\end{abstract}

\section{The problem}

The problem of the many was posed by \citet{Unger1980} and, in the version with a cat and its hairs, by \citet{Geach1980}. Consider a cloud. Around its edge are droplets that are neither clearly part of it nor clearly not. For each such droplet, there is an aggregate that includes it and one that does not, and the two aggregates differ only by that droplet. There are therefore very many aggregates, each with as good a claim as any other to be the cloud. Unger drew the conclusion that there are no clouds, and no ordinary objects of any kind. Others have accepted that there are many clouds where we see one and sought to explain why we say \enquote{one}: \citet{Lewis1993} proposed that the many clouds are almost identical, sharing nearly all their parts, and that ordinary counting is counting by almost-identity. \citet{Lowe1982} and \citet{Geach1980} appealed to sortal-relative counting. \citet{McGee2000} gave a supervaluationist treatment on which \enquote{there is exactly one cloud} is determinately true although no candidate is determinately the cloud. Lewis's paper offered the supervaluationist solution first and the almost-identity solution as a supplement for those who reject it. \citet{Weatherson2003} argues that the problem arises far more widely than for clouds and cats.

We defend the answer \enquote{one}. The candidates are not each a cloud, and the vagueness of \enquote{cloud} selects one of them by supervaluation.

\section{The argument}

\begin{definition}\label{d:cand}
A \emph{candidate} in a region is an aggregate of droplets that every admissible precisification of \enquote{cloud} could count as a cloud: it contains all the droplets that are clearly part of the cloud and none that are clearly not. A \emph{precisification} is an admissible sharpening of \enquote{cloud}, one that respects its clear cases and its penumbral connections \citep{Fine1975}.
\end{definition}

\begin{principle}[Maximality]\label{p:max}
\enquote{Cloud} is a maximal sortal: on every precisification, no two distinct clouds share most of their droplets. In particular, a large proper part of a cloud is not itself a cloud.
\end{principle}

\ref{p:max} is a penumbral connection of the kind \citet{Fine1975} identifies. Many sortals for ordinary objects are maximal in this sense \citep{Sider2003}: a house with one brick removed from its wall is not a further house inside the first, and two massively overlapping aggregates of bricks are not two houses. \enquote{Cloud} is maximal in the same way. The principle is a constraint on each precisification, not a claim that any particular candidate is the cloud.

\begin{proposition}\label{p:one}
On every precisification, there is exactly one cloud in the region.
\end{proposition}

\begin{proof}
A precisification draws a sharp boundary for \enquote{cloud}. At least one candidate is a cloud on it, since the region is a clear case of containing a cloud and every precisification respects clear cases. Suppose two distinct candidates $c_1$ and $c_2$ were clouds on the same precisification. Both contain all the droplets that are clearly part of the cloud, and these are most of the droplets of each, since the candidates differ only at the edge. By \ref{p:max}, $c_1 = c_2$, a contradiction. So exactly one candidate is a cloud on each precisification.
\end{proof}

\begin{corollary}\label{c:super}
\enquote{There is exactly one cloud in the region} is supertrue, while for no candidate $c$ is \enquote{$c$ is the cloud} supertrue.
\end{corollary}

\begin{proof}
The first part is Proposition~\ref{p:one}. For the second, each candidate is, by Definition~\ref{d:cand}, not counted as a cloud by some precisification that counts another candidate instead.
\end{proof}

Corollary~\ref{c:super} answers both horns of Unger's dilemma. There are not many clouds, since the counting sentence is true on every precisification. There is not no cloud, for the same reason. And no arbitrary choice has been made, since the semantics does not single out any candidate.

\section{Supervaluation and the epistemic account}

MET-18 argues that vague predicates have sharp but unknowable boundaries: \enquote{cloud} has a single extension, fixed by use, and does not merely have many admissible precisifications. The appeal to supervaluation here does not conflict with that view, and we state how the two fit together.

On the epistemic account, the actual pattern of use fixes one precisification, the \emph{intended} one; truth is truth on it. The other admissible precisifications are those the intended one cannot be discriminated from by the speakers, in the sense of MET-18's margin model. A sentence true on all admissible precisifications is therefore true on every interpretation indiscriminable from the actual one, and so known, in the sense of MET-18-a. Supertruth, so understood, is not a definition of truth but a model of determinate, that is knowable, truth. Read in this way, Corollary~\ref{c:super} says two things. It is known that there is exactly one cloud in the region. And of no candidate is it known that it is the cloud. The epistemic account adds what the supervaluationist reading leaves out: that exactly one candidate is in fact the cloud, namely the one counted by the intended precisification, and that its identity is unknowable for the same reason that the location of the boundary of \enquote{bald} is unknowable.

The two descriptions therefore differ only on whether \enquote{$c$ is the cloud} has a truth value for a given candidate $c$. On the supervaluationist definition of truth it does not; on the epistemic account it does, and it is true for exactly one $c$. The conclusion of the present manuscript, that there is one cloud and that the vagueness of \enquote{cloud} selects it, holds on both. The selection is effected by use, its existence is established by supervaluation over the indiscriminable precisifications, and its result is not knowable.

\section{Objections}

\begin{objection}[Many, but almost one]
Each candidate has everything it takes to be a cloud: it is a large, visible, dense aggregate of droplets. If one of them is a cloud, they all are, and the answer is that there are many clouds, almost identical, which we count as one \citep{Lewis1993}.
\end{objection}

\begin{reply}
The candidates are intrinsically alike in the respects that matter for being cloud-like, but being a cloud is not intrinsic: by \ref{p:max}, it depends on not overlapping massively with another cloud. Counting by almost-identity, moreover, is not counting. The sentence \enquote{there is exactly one cloud} is about identity, and on Lewis's supplementary view it is false, while we say it is true. A view that keeps it true is preferable if it is available, and Proposition~\ref{p:one} shows that it is.
\end{reply}

\begin{objection}[Nihilism]
\citet{Unger1980} argues that if no candidate is privileged, none is a cloud, and there are no clouds.
\end{objection}

\begin{reply}
The inference from \enquote{no candidate is privileged} to \enquote{no candidate is a cloud} is invalid. On every precisification one candidate is privileged, and on the epistemic account one is privileged in fact, by use. What is true is that none is knowably privileged, and from that nothing follows about whether there are clouds.
\end{reply}

\begin{objection}[Composition]
MET-06 restricts composition to things that share a common fate. If the restriction left only one aggregate of droplets in the region, the problem would not arise; if it leaves several, it seems to yield many clouds.
\end{objection}

\begin{reply}
It leaves several. As MET-06-a observes, the pluralities that differ by a few droplets at the edge share the common fate of the mass, so each candidate aggregate exists as a composite; what the criterion excludes are pluralities that include droplets carried off by a different current, which are therefore not candidates. Composition says which aggregates there are, and leaves to the semantics of \enquote{cloud} which of them is a cloud. The many composites are not many clouds: by Proposition~\ref{p:one}, only one of them is a cloud.
\end{reply}

\section{Conclusion}

There is one cloud. The aggregates that differ at the edge exist, and are candidates for being the cloud, but on every admissible precisification exactly one is a cloud, because \enquote{cloud} is a maximal sortal. The vagueness of \enquote{cloud} selects one candidate on each precisification, and the counting sentence is supertrue. On the epistemic account of MET-18, supertruth models what is knowable: it is known that there is one cloud, and exactly one candidate is the cloud, though no one can know which.

\begin{thebibliography}{99}\small
\bibitem[Fine(1975)]{Fine1975} Fine, K. (1975). Vagueness, truth and logic. \emph{Synthese}, 30, 265--300.
\bibitem[Geach(1980)]{Geach1980} Geach, P.~T. (1980). \emph{Reference and Generality}, 3rd edition. Cornell University Press.
\bibitem[Lewis(1993)]{Lewis1993} Lewis, D. (1993). Many, but almost one. In J.~Bacon, K.~Campbell and L.~Reinhardt (eds.), \emph{Ontology, Causality and Mind: Essays in Honour of D.~M. Armstrong}. Cambridge University Press.
\bibitem[Lowe(1982)]{Lowe1982} Lowe, E.~J. (1982). The paradox of the 1,001 cats. \emph{Analysis}, 42(1), 27--30.
\bibitem[McGee and McLaughlin(2000)]{McGee2000} McGee, V. and McLaughlin, B. (2000). The lessons of the many. \emph{Philosophical Topics}, 28(1), 129--151.
\bibitem[Sider(2003)]{Sider2003} Sider, T. (2003). Maximality and microphysical supervenience. \emph{Philosophy and Phenomenological Research}, 66(1), 139--149.
\bibitem[Unger(1980)]{Unger1980} Unger, P. (1980). The problem of the many. \emph{Midwest Studies in Philosophy}, 5, 411--467.
\bibitem[Weatherson(2003)]{Weatherson2003} Weatherson, B. (2003). Many many problems. \emph{The Philosophical Quarterly}, 53(213), 481--501.
\end{thebibliography}

\end{document}
